\section{Resultados e Testes}

Esta seção mostra como se verificou que o validador faz o que a
especificação promete.

\subsection{Casos de teste}

A Tabela~\ref{tab:testes} reúne os casos usados. Eles cobrem os
dois padrões, valores extremos e os erros mais comuns de
digitação. O programa de testes automatizados executa cada caso e
compara a resposta com a coluna ``Esperado''.

\begin{table}[ht]
\centering
\small
\renewcommand{\arraystretch}{1.1}
\begin{tabular}{llp{6.3cm}}
\toprule
Entrada & Esperado & Motivo \\
\midrule
\texttt{ABC1234}   & Válida   & padrão anterior, sem hífen \\
\texttt{ABC-1234}  & Válida   & padrão anterior, com hífen \\
\texttt{AAA0000}   & Válida   & padrão anterior, valores extremos \\
\texttt{BRA2E19}   & Válida   & Mercosul \\
\texttt{QTP9Z01}   & Válida   & Mercosul, letra de A a Z na 5ª posição \\
\texttt{ABC1C34}   & Válida   & Mercosul convertido de \texttt{ABC1234} \\
\texttt{ZZZ9Z99}   & Válida   & Mercosul, valores extremos \\
\midrule
\texttt{ABC-1D23}  & Inválida & hífen no Mercosul \\
\texttt{AB12C34}   & Inválida & apenas 2 letras iniciais \\
\texttt{ABC12D3}   & Inválida & letra na 6ª posição \\
\texttt{BRA2E1}    & Inválida & apenas 6 caracteres \\
\texttt{ABC12345}  & Inválida & algarismo a mais \\
\texttt{ABC123}    & Inválida & algarismo a menos \\
\texttt{ABCD123}   & Inválida & quatro letras seguidas \\
\texttt{ABC1DD3}   & Inválida & letra no lugar de algarismo \\
\texttt{ABC--1234} & Inválida & dois hífens \\
\texttt{ABC\textvisiblespace 1234} & Inválida & espaço no meio \\
\texttt{abc1234}   & Inválida & letras minúsculas \\
\texttt{ÁBC1234}   & Inválida & letra com acento \\
(vazia)            & Inválida & cadeia vazia \\
\bottomrule
\end{tabular}
\caption{Casos de teste do validador.}
\label{tab:testes}
\end{table}

\subsection{Verificação exaustiva}

Testar casos escolhidos à mão pode deixar passar algum erro. Por
isso o validador também foi comparado com a expressão regular
$\Lt^3(\Hf?\,\Nm^4 \mid \Nm\,\Lt\,\Nm^2)$ em \emph{todas} as
cadeias possíveis sobre as classes $\{\Lt, \Nm, \Hf\}$ de
comprimento até 9, o que dá
$\sum_{n=0}^{9} 3^n = 29\,524$ cadeias.

Esse conjunto basta por dois motivos. Primeiro, símbolos da mesma
classe se comportam do mesmo modo, então um representante por
classe cobre todos. Segundo, nenhuma placa aceita tem mais de 8
caracteres, e os estados finais só saem para o estado de erro;
logo toda cadeia mais longa é rejeitada. Caracteres fora do
alfabeto também são rejeitados, pois não têm transição definida.

\subsection{Quantas placas o validador aceita}

\begin{table}[ht]
\centering
\renewcommand{\arraystretch}{1.2}
\begin{tabular}{llr}
\toprule
Formato & Cálculo & Cadeias aceitas \\
\midrule
Anterior, sem hífen & $26^3 \cdot 10^4$ & $175\,760\,000$ \\
Anterior, com hífen & $26^3 \cdot 10^4$ & $175\,760\,000$ \\
Mercosul & $26^3 \cdot 10 \cdot 26 \cdot 10^2$ & $456\,976\,000$ \\
\midrule
Total de cadeias & & $808\,496\,000$ \\
\bottomrule
\end{tabular}
\caption{Quantidade de cadeias aceitas.}
\label{tab:contagem}
\end{table}

As grafias com e sem hífen representam a mesma placa. Contando cada
placa uma única vez, o formato admite
$175\,760\,000 + 456\,976\,000 = 632\,736\,000$ placas distintas.
Esse é o tamanho do espaço de combinações possíveis, e não o número
de placas realmente emitidas, que é menor e depende dos cadastros
oficiais.